Enabling IMA and EVM can be complicated because of the need to ensure trustworthiness of the installed IMA policy and signatures.
We recommend at least reading the IBM IMA how-to documentation \cite{IBM:IMA-utilities-evmctl}, sections ``IMA CA Key and Certificate'' to ``Sign and Install a Custom Policy''.

On Fedora, dracut is the tool used to build \initramfs.
Dracut has a standard module for loading IMA/EVM keys and enabling IMA/EVM during boot named \verb|integrity|.
It is necessary to remove the \verb|ima-policy-load.sh| script from this module, since Fedora's SELinux policy is not designed to load during boot and our IMA policy relies on the SELinux policy.
The \verb|integrity| module loads IMA keys from \verb|/etc/keys/ima/| as described in \cite{IBM:IMA-utilities-evmctl}.
Before that, it loads the EVM keys from \verb|/etc/keys/evm/|;
thus, another key with a similar setup (or the same key) as generated for IMA in the ``IMA Signing Key and Certificate'' section must also be placed there to serve as the EVM signature verification key.

As discussed briefly in \zcref[S]{sec:implementation-details}, sometimes it is necessary to enable EVM HMACs as well.
This requires creating a kernel masterkey (we make a trusted masterkey since we are using a TPM) and a child encrypted EVM HMAC key.
\begin{lstlisting}[style=bash]
tpm2_createprimary --key-algorithm=rsa2048 --key-context=key.ctxt
tpm2_evictcontrol -c key.ctxt 0x81000001
KMK=$(keyctl add trusted kmk-trusted "new 32 keyhandle=0x81000001" @u)
keyctl pipe $KMK > /etc/keys/kmk-trusted.blob
EVM=$(keyctl add encrypted evm-key "new trusted:kmk-trusted 64" @u)
keyctl pipe $EVM > /etc/keys/evm-trusted.blob
\end{lstlisting}
% TODO : cite? https://web.archive.org/web/20250709141759/https://docs.redhat.com/en/documentation/red_hat_enterprise_linux/9/html/managing_monitoring_and_updating_the_kernel/enhancing-security-with-the-kernel-integrity-subsystem_managing-monitoring-and-updating-the-kernel#working-with-trusted-keys_enhancing-security-with-the-kernel-integrity-subsystem